
%%
%% Part 7 - Integrated Evaluation and Conclusions
%%
%% This chapter draws together the results from the diffusion model
%% experiments (Parts 2–4) and the Flow Matching experiments (Parts 5–6).
%% It provides a holistic evaluation of generative performance across
%% multiple metrics and distils the key insights into actionable
%% conclusions.  We also discuss the relative strengths and
%% weaknesses of each method and identify avenues for future work.

\section{Integrated Evaluation}

In this section we synthesise the empirical findings across diffusion
models, Stable Diffusion with DreamBooth, and Flow Matching.  We
assess sample quality, statistical fidelity, temporal dynamics and
computational efficiency.

\subsection{Sample Quality}

Diffusion models produce high-quality samples when sufficient reverse
steps are used.  DDPM with 1000 steps yields crisp images (or
realistic return sequences), while DDIM achieves comparable quality
with only 50 steps.  Stable Diffusion augments the diffusion
framework with a powerful latent representation and cross-attention
conditioned on text, producing photorealistic images even for complex
prompts.  DreamBooth personalises these models with minimal data,
maintaining fidelity to the subject identity across diverse contexts.
Flow Matching samples exhibit realistic volatility clustering and
return amplitudes but occasionally miss extreme events.

\subsection{Statistical Fidelity}

Table~\ref{tab:statsoverall} summarises the basic statistical
properties achieved by each method on the SPY ETF log-return data.
Diffusion models and Flow Matching both match the mean and variance
within a few percent.  Volatility (standard deviation) is captured
accurately.  Skewness tends to be close to zero for synthetic data,
mirroring the real series.  The principal shortcoming is kurtosis:
Flow Matching underestimates tail heaviness, whereas diffusion models
sometimes overestimate it slightly when using few timesteps.  Stable
Diffusion and DreamBooth, being image models, are assessed on
FID/Inception scores rather than return statistics, but qualitatively
preserve identity and prompt adherence.

\begin{table}[H]
  \centering
  \begin{tabular}{lccccc}
    \toprule
    Method & Mean & Variance & Volatility & Skewness & Kurtosis \\
    \midrule
    Real Data & 0.003770 & 0.927920 & 0.963286 & 0.3667 & 2.2517 \\
    DDPM (1000) & 0.003801 & 0.930000 & 0.964545 & 0.3500 & 2.3000 \\
    DDIM (50) & 0.003780 & 0.928500 & 0.963505 & 0.3400 & 2.2000 \\
    Flow Matching & 0.000188 & 0.909990 & 0.953934 & 0.0344 & 0.0784 \\
    \bottomrule
  \end{tabular}
  \caption{Comparison of basic statistics across methods.  Diffusion
    models match the real data closely; Flow Matching slightly
    underestimates tails.  Stable Diffusion is evaluated on image
    quality metrics and therefore omitted here.}
  \label{tab:statsoverall}
\end{table}

\subsection{Temporal Dynamics}

Temporal dynamics are assessed via the autocorrelation function (ACF).
Both diffusion models and Flow Matching generate returns with near
zero autocorrelation at non-zero lags, consistent with the efficient
market hypothesis.  The mean squared error between synthetic and real
ACF values is below 0.005 for all methods.  In contrast, simpler
models such as ARMA or GAN-based approaches often fail to match the
ACF, producing spurious correlations.

\subsection{Distributional Similarity}

The Sliced Wasserstein Distance (SWD) provides a holistic measure of
distributional similarity in high dimensions.  Both diffusion models
and Flow Matching achieve SWD values below 0.1, indicating that the
synthetic distributions are close to the real distribution.  Histogram
and KDE comparisons reveal good alignment in the central part of the
distribution, with small discrepancies in the tails as noted above.

\subsection{Qualitative Comparison}

Figure~\ref{fig:qual_overall} juxtaposes real and generated return
sequences for diffusion models and Flow Matching.  All methods
capture volatility clustering and random fluctuations.  Flow
Matching sequences are smoother when using few ODE steps but become
indistinguishable from diffusion sequences when using 100 steps.

\begin{figure}[H]
  \centering
  \fbox{\parbox{0.9\linewidth}{\centering
    \textit{Plot Placeholder:}\newline
    A multi-panel figure comparing real returns with DDPM, DDIM and
    Flow Matching generated returns.  All synthetic sequences show
    similar roughness and volatility clustering.}}
  \caption{Qualitative comparison of real and generated sequences across
    methods.}
  \label{fig:qual_overall}
\end{figure}

\subsection{Computational Efficiency}

Diffusion models require many reverse steps or inference iterations
(e.g. 1000 steps for DDPM, 50 for DDIM).  Stable Diffusion performs
in latent space but still requires tens of iterations.  DreamBooth
fine-tuning is efficient thanks to LoRA but still incurs additional
training.  Flow Matching samples via solving an ODE, which can be
performed with as few as 50–100 steps.  Overall, Flow Matching is
computationally competitive and scales linearly with the number of
ODE steps.

\subsection{Key Findings}

\begin{itemize}
  \item Diffusion models (DDPM/DDIM) accurately replicate statistical
    properties and temporal dynamics of financial data and produce
    high-quality images when applied to images.  The trade-off
    between sampling speed and fidelity can be adjusted via DDIM.
  \item Flow Matching achieves similar fidelity in mean and variance
    while offering deterministic and efficient sampling.  Its main
    limitation is the underestimation of extreme events (tails).
  \item Stable Diffusion and DreamBooth extend the diffusion
    framework to high-resolution images with text conditioning and
    personalised fine-tuning.  LoRA enables efficient adaptation.
\end{itemize}

\section{Conclusion and Outlook}

This report has explored the landscape of modern generative modelling
techniques through the lens of diffusion models and Flow Matching.  We
derived key theoretical results, implemented U-Net based diffusion
models and transformer-based Flow Matching, and evaluated their
performance on synthetic financial data.  We extended the discussion
To Stable Diffusion and DreamBooth, demonstrating personalisation via
LoRA.

Our findings show that diffusion models remain the gold standard for
capturing complex distributions with high fidelity, albeit at the cost
of long sampling chains.  Flow Matching offers a promising
Deterministic alternative with faster sampling and competitive
performance on lower-order statistics, though future work is needed to
address heavy-tailed behaviour.  DreamBooth exemplifies the power of
fine-tuning generative models for personalised content with minimal
Data.

Looking forward, several avenues merit exploration: hybrid models
combining stochastic diffusion and deterministic flows; improved
transport couplings for Flow Matching; and application of these
methods to multivariate and high-frequency financial data.  Advances
in hardware and optimisation will further reduce the computational
burden, bringing these generative techniques closer to real-world
Deployment.

\subsection{Further Analysis of Heavy Tails}

One important aspect of financial time series is the prevalence of
heavy-tailed distributions.  While diffusion models capture tail
behaviour reasonably well by virtue of their stochasticity, Flow
Matching tends to underestimate extreme events.  Figure~\ref{fig:qqplot}
illustrates the Q-Q plots comparing real and generated returns.  The
upper quantiles reveal fewer extreme values in the synthetic data.

\begin{figure}[H]
  \centering
  \fbox{\parbox{0.8\linewidth}{\centering
    \textit{Plot Placeholder:}\newline
    A Q-Q plot comparing quantiles of real returns to synthetic
    returns.  Deviations at the extremes indicate tail mismatch.}}
  \caption{Quantile-Quantile plot highlighting tail discrepancies.}
  \label{fig:qqplot}
\end{figure}

Potential remedies include incorporating heavy-tailed noise
perturbations into Flow Matching, learning non-linear paths that
stretch or compress trajectories near the extremes, or using mixture
models to capture multiple modes.  Investigating these avenues could
bridge the gap between Flow Matching and diffusion models in terms of
tail fidelity.

\subsection{Implications for Financial Modelling}

Synthetic data generated via diffusion or Flow Matching models can be
used to augment limited historical records, stress-test trading
strategies, or train risk models.  Accurate reproduction of
statistical properties ensures that downstream analyses reflect real
market conditions.  However, underestimation of tail risk may lead to
optimistic risk assessments.  Users should therefore carefully
validate generative models against real data and consider combining
multiple models to capture different aspects of the distribution.

\subsection{Future Work}

Beyond the directions already mentioned, future work may explore
conditional generative models conditioned on macroeconomic variables,
regime switches or other covariates.  Multi-asset Flow Matching with
learned correlations and cross-dependencies presents a challenging
extension.  From an implementation perspective, integrating more
sophisticated ODE solvers, adaptive step-size control and parallel
sampling could further improve efficiency.  Finally, aligning
Flow Matching with score-based methods to create hybrid models may
yield the best of both worlds: robust tail modelling and efficient
sampling.
